\documentclass[11pt]{amsart}
\usepackage{calc,amssymb,amsmath,ulem,stmaryrd,fullpage}
\usepackage{enumerate}
\usepackage{alltt}
\RequirePackage[dvipsnames,usenames]{color}
\let\mffont=\sf
\normalem
%\input{kmacros3.sty}
\input{xy}
\xyoption{all}
\usepackage{hyperref}
\usepackage{graphicx}
\usepackage[all,cmtip]{xy}
\usepackage{verbatim}
\usepackage[left=0.95in,top=0.95in,right=0.95in,bottom=0.95in]{geometry}
\newcommand{\tauCohomology}{T}
\newcommand{\FCohomology}{S}
\usepackage{mathrsfs}


\theoremstyle{theorem}
%\newtheorem*{mainthm*}{Main Theorem}

\newcommand{\note}[1]{\marginpar{\sffamily\tiny #1}}

\def\todo#1{\textcolor{Mahogany}%
{\sffamily\footnotesize\newline{\color{Mahogany}\fbox{\parbox{\textwidth-15pt}{#1}}}\newline}}

\renewcommand{\O}{\mathcal O}
\newcommand{\ie}{{\itshape i.e.} }
\newcommand{\roundup}[1]{\lceil #1 \rceil}
\newcommand{\rounddown}[1]{\lfloor #1 \rfloor}
\renewcommand{\to}[1][]{\xrightarrow{\ #1\ }}
\newcommand{\onto}[1][]{\protect{\xrightarrow{\ #1\ }\hspace{-0.8em}\rightarrow}}
\newcommand{\into}[1][]{\lhook \joinrel \xrightarrow{\ #1\ }}
%\newcommand{DBDual}[1]{{\underline{\omega}_{#1}}}
\newcommand{\DBDual}[1]{{{\uuline \omega}{}^{\mydot}_{#1}}}
\newcommand{\Tt}{{\mathfrak{T}}}
%\newcommand{\bn}{{\mathfrak{n}} }
\newcommand{\sol}[1]{ \vskip 6pt{\color{blue} {\bf Solution:} #1} \vskip 6pt}

\newcommand{\bR}{{\mathbb{R}}}
\newcommand{\va}{{\vec{a}}}
\newcommand{\vb}{{\vec{b}}}
\newcommand{\vzero}{{\vec{0}}}
\newcommand{\vi}{{\vec{i}}}
\newcommand{\vj}{{\vec{j}}}
\newcommand{\vk}{{\vec{k}}}
\newcommand{\vh}{{\vec{h}}}
\newcommand{\vr}{{\vec{r}}}
\newcommand{\vu}{{\vec{u}}}
\newcommand{\vv}{{\vec{v}}}
\newcommand{\vw}{{\vec{w}}}
\newcommand{\vx}{{\vec{x}}}
\newcommand{\vy}{{\vec{y}}}
\newcommand{\vz}{{\vec{z}}}
\newcommand{\vT}{{\vec{T}}}
\newcommand{\vN}{{\vec{N}}}
\newcommand{\vF}{{\vec{F}}}
\newcommand{\vG}{{\vec{G}}}
\newcommand{\bF}{\mathbb{F}}
\newcommand{\bQ}{\mathbb{Q}}
\newcommand{\bZ}{\mathbb{Z}}
\newcommand{\bP}{\mathbb{P}}
\newcommand{\bA}{\mathbb{A}}
%\newcommand{\cF}{\mathcal{F}}
%\newcommand{\cG}{\mathcal{G}}
%\newcommand{\cH}{\mathcal{H}}
\newcommand{\image}{\mathrm{im}\,}
\newcommand{\coker}{\mathrm{coker}\,}
\newcommand{\Hom}{\mathrm{Hom}\,}
\newcommand{\Supp}{\mathrm{Supp}\,}
\newcommand{\Ann}{\mathrm{Ann}\,}
\newcommand{\sH}{\mathscr{H}\,}
\newcommand{\cI}{\mathcal{I}}
\newcommand{\cJ}{\mathcal{J}}
\newcommand{\cE}{\mathcal{E}}
\newcommand{\cF}{\mathcal{F}}
\newcommand{\cG}{\mathcal{G}}
\newcommand{\cH}{\mathcal{H}}
\newcommand{\cO}{\mathcal{O}}
\newcommand{\cM}{\mathcal{M}}
\newcommand{\cN}{\mathcal{N}}
\newcommand{\cK}{\mathcal{K}}
\newcommand{\fram}{\mathfrak{m}}
\newcommand{\Div}{\mathrm{Div}}

\begin{document}
\title{Worksheet \#4 -- Math 6140\\Spring 2019}
\author{Due Friday, March 8th}

\maketitle
You may work in groups of up to 3.  Only one worksheet needs to be turned in per group.
\vskip 9pt
\noindent
We begin with some definitions.
\vskip 9pt
{\bf Definition:} {\it 
A Cartier divisor $D$ on a normal projective variety $X$ is called:
\begin{itemize} 
\item[(a)] \emph{very ample} if $|D|$ induces an embedding into projective space.
\item[(b)] \emph{ample} if $nD$ is very ample for some $n > 0$.
\item[(c)] \emph{semi-ample} if $|nD|$ is base point free for some $n > 0$.
\item[(d)] \emph{big} if the induced rational map to projective space $\phi_{|nD|} : X \dashrightarrow \bP^m$ is birational onto its image.
\end{itemize}
}
\vskip 9pt
The more standard definition of ample is the following:
{\it
\begin{itemize}
\item[(b')]  A Cartier divisor $D$ is \emph{ample} if for any coherent sheaf $\cF$, there exists an integer $n_0$ so that $\cF \otimes \O_X(nD) =: \cF(nD)$ is globally generated for all $n \geq n_0$.  
\end{itemize}
}
These are indeed equivalent, as we'll show shortly.  
\vskip 9pt
\noindent
{\bf 1.}  Suppose that $D$ is very ample and that $E$ is such that $|E|$ is base point free.  Prove that $D+E$ is also very ample.  In particular, conclude that if $D$ is very ample, so is $2D$, $3D$, etc, and likewise with ample.  
\vskip 3pt
\emph{Hint:} You have maps $\phi_{|D|} : X \to \bP^r$ and $\phi_{|E|} : X \to \bP^s$.  Consider the induced product map $X \to \bP^r \times \bP^s$.  Show that this corresponds to a linear system that is contained within the linear system $|D+E|$.
\newpage
\noindent
{\bf 2.}  Suppose that $D$ is very ample.  Prove that $D$ satisfies the condition (b') above.  Conclude that ample divisors in the sense of (b) are ample in the sense of (b').
\vskip 3pt
\emph{Hint:} Suppose that $i : X \to \bP^n$ is the induced closed embedding.  Prove that $\Gamma(X, \cF(nD)) = \Gamma(\bP^n, (i_* \cF) \otimes_{\O_{\bP^n}} \O_{\bP^n}(n))$.
\vskip 250pt
\noindent
{\bf 3.}  Suppose that $D$ is a semi-ample divisor on $X$.  Further suppose that there are two points $p$ and $q$ so that if $D' \sim nD$ and $p$ or $q$ are in the support of $D'$, then $p$ and $q$ are in the support of $D'$.  Prove that $D$ cannot be ample.  
\vskip 3pt
\emph{Hint} Use Lemma 13.28 from the text.
\newpage
The following can also be found in Cutkosky/Hartshorne.  Suppose $\phi : X \to \bP^n$ corresponds to a vector space $V \subseteq \Gamma(X, \O_X(D))$ where $D$ is a Cartier divisor.  
\begin{itemize}
\item[(i)]  $\phi$ is regular at a point $p$ if and only if the canonical map $V \to \O_X(D) / I_p(D) \cong k$ surjects.  (This map is just the restriction map, followed by quotienting).
\item[(ii)]  $\phi$ is injective and regular if and only if for every points $p \neq q \in X$, we have that the natural map $V \to (\O_X(D) / I_p(D)) \oplus (\O_X(D) / I_q(D)) \cong k^2$ surjects.  
\item[(iii)]  $\phi$ is a closed embedding if and only if it is injective and regular and for each $p \in X$, we have that the set $\{ s_p \in \fram_p \O_X(D)_p | s \in V \}$ spans the dual tangent space $(\fram_p \O_X(D)) / (\fram_p^2 \O_X(D)) \cong \fram_p/\fram_p^2$.
\end{itemize}
Condition (ii) is usually called $|V|$ separating points and (iii) is called $|V|$ separating tangent vectors.  We've basically done (i) and (ii), we will work on showing (iii) now.  
\vskip 9pt
\noindent
{\bf 4.}  Consider $V = \Gamma(X, \O_X(D))$.  Prove that the condition from (iii) holds if $D$ is very ample in the sense of (b). 
\vskip 3pt
\emph{Hint:} Observe that $\O_{\bP^n}(1)$ satisfies the condition in (iii).  
\vskip 170pt
\noindent
{\bf 5.}  Prove that the condition from (iii) really does imply $\phi$ is a closed embedding.  
\vskip 3pt
\emph{Hint:}  You may use the following algebra fact.  If $f : R \to S$ is a local map of local rings such that $S$ is a finite $R$-module, $R/\fram_R \to S/\fram_S$ is an isomorphism, and $\fram_R/\fram_R^2 \to \fram_S/\fram_S^2$ is surjective, then $f$ is surjective.  
\newpage
\noindent
{\bf 6.}  Show that (b') implies (b) and thus our different characterizations of ampleness all coincide.
\vskip 140pt
\noindent
{\bf 7.}  Suppose that $i : Y \hookrightarrow  X$ is a closed embedding of projective varieties and that $D \subseteq X$ is an ample (resp. semi-ample) divisor such that $Y \not\subseteq \Supp D$.  Prove that $i^* D = D|_Y$ is ample (resp. semi-ample).
\vskip 140pt
\noindent
{\bf 8.}  Give an example of a projective variety $X$, a closed subvariety $Y \subseteq X$, and a big divisor $D$ on $X$ such that $Y \not\subseteq \Supp D$, so that $D|_Y$ is not big. 
\vskip 140pt
\noindent
{\bf 8.}  Give an example of a surjective map of projective varieties (of dimension $> 0$) $\phi : Y \to X$ and a big divisor $D$ on $X$ such that $\phi^* D$ is not big.  
\vskip 3pt
\emph{Hint:} If $D$ is big, then the map associated to $|nD|$ is injective at most points.
\end{document}

